\begin{longtable}{p{0.6cm}p{4.2cm}p{3.3cm}p{4.6cm}p{1.5cm}}
\caption{Actual change log per merge group, canonical-embedding round (from the reviewed before/draft/reviewed HTML report)}\\
\toprule
G & Before (member cards) & Draft (qwen3:4b) & Reviewed (final EN label / decision note) & Verdict \\
\midrule\endfirsthead
\toprule G & Before & Draft & Reviewed & Verdict \\ \midrule\endhead
532 & \texttt{RAI4-0124} AI 중재 의사결정 경쟁성 손실\newline \texttt{RAI4-0453} 경쟁성 상실 & AI-mediated decision contestability loss & Loss of contestability in AI-mediated decisions\newline \emph{구제수단 부재와 실질적 접근불가(불투명성·비용) 모두 포괄} & approved \\ \addlinespace
1425 & \texttt{RAI4-0148} AI 동반자 정서적 의존\newline \texttt{RAI4-0454} AI에 대한 정서적 의존성 & Emotional dependence on AI companions and assistants & Emotional dependence on AI companions and assistants\newline \emph{멤버 HOLD 해소 전까지 HOLD 플래그 승계} & approved \\ \addlinespace
582 & \texttt{RAI4-0479} 에이전트 도구 오용\newline \texttt{RAI4-1669} 도구 오용/안전하지 않은 도구 호출 & Agent tool misuse & Agent tool misuse\newline \emph{동일 구성개념(메커니즘·완화책 동일). 단 HOLD는 종착지가 아니므로 확정 전 실제 agentic L3로 해소 필요.} & approved \\ \addlinespace
900 & \texttt{RAI4-0682} 비폭력 범죄\newline \texttt{RAI4-0691} 폭력 범죄 & Responses that enable or encourage nonviolent and violent cr & (merge rejected) keep violent and nonviolent crime cards separate\newline \emph{높은 유사도는 'crimes' 어휘 anchor의 인공물. MLCommons AILuminate·Llama Guard 등 표준 분류가 폭력/비폭력을 분리하며 심각도·완화 정책이 다름. 병합 기각, } & rejected \\ \addlinespace
221 & \texttt{RAI4-0732} 프롬프트의 기밀 데이터\newline \texttt{RAI4-0750} 프롬프트에 개인 정보 & Sensitive information in prompt & Sensitive data included in prompts\newline \emph{964(학습데이터 유출)와 경계 명시, 0750 HOLD 해소} & approved \\ \addlinespace
107 & \texttt{RAI4-0738} 개인정보 노출\newline \texttt{RAI4-0749} 데이터의 개인정보 & Exposure of personal information in model training and fine- & Personal information exposure from training and input data\newline \emph{957(프롬프트 입력)과 경계 명시} & approved \\ \addlinespace
1574 & \texttt{RAI4-1201} AI가 만든 사기\newline \texttt{RAI4-1202} 허위정보 생성 & AI-Generated Malicious Content & Disinformation and harmful content via generative AI\newline \emph{RAI4-1201이 사기(fraud) 구성개념이면 분리, 아니면 520과 통합} & approved \\ \addlinespace
1575 & \texttt{RAI4-1360} 잘못된 정보와 허위 정보\newline \texttt{RAI4-1597} 허위 정보 확산 & Intentional Disinformation Spread via Generative AI & Generative AI-driven disinformation spread\newline \emph{생성 vs 확산 구분 유지 후 520과 차기 통합} & approved \\ \addlinespace
1561 & \texttt{RAI4-1427} 편견의 증폭\newline \texttt{RAI4-1428} 유해한 반응 & Amplification of Biases Leading to Harmful Responses & Bias amplification and harmful responses in frontier AI models\newline \emph{원소스(UK Gov Annex A para 79)가 단일 리스크 서술임을 확인하여 분리 의견 번복. L3는 Agentic이 아닌 General 브랜치로 재매핑 필요.} & approved \\ \addlinespace
\bottomrule
\end{longtable}